\section{国际法律责任}
\begin{enumerate}
  \item 主体：
        \begin{enumerate}
          \item 主权国家：最最主要的国际法主体。
          \item \zd{政府间}国际组织：\zd{联合国，世贸}。政府间和非政府间（国际红十字会）：\qd{依政府间协议}成立和活动。
          \item 某些特定的民族解放组织：巴解（巴勒斯坦民族解放组织）。不会考。
        \end{enumerate}
        \begin{kk}{国家的构成要素和基本权利}
          \begin{itemize}
            \item 不是直接考点但要了解。
            \item 构成要素：定居的居民、确定的领土、政府、主权。
            \item 基本权利：独立平等自保管辖（属人属地保护普遍性）。
          \end{itemize}
        \end{kk}
  \item 考点，与主权国家有关的三大制度：
        \begin{enumerate}
          \item 豁免：
                \begin{enumerate}
                  \item 绝对豁免：其他国家管不着。例外：自愿放弃，特定案件，\zd{明示默示}。
                  \item \zd{诉行为}：\qd{起诉、应诉、反诉、作为利害关系人介入诉讼}。\zd{出庭 \bdy 放弃}。
                  \item 判决不履行 \yjt 没有执行的效力（\qd{国家的财产不得被查封冻结扣押}）。\zd{管辖豁免的放弃和执行豁免的放弃是分开的}。
                  \item 限制豁免（相对豁免）。国商。未生效，理论状态。
                \end{enumerate}
          \item 承认（重要考点）。
                \begin{enumerate}
                  \item 国家和政府间国际组织有资格。
                  \item 明示默示（外交关系，互设使馆；领事关系；政治性条约（双边税收协定）；\zd{联合国正式投票}）。
                  \item \zd{新国家（领土变更）、新政府（革命、内战）}。考点。
                  \item 单方行为（最最重要）。\qd{拒绝承认不是剥夺人家国际法主体资格的方式}。
                  \item 承认新国家\zd{不可逆}；承认新政府意味着\zd{对旧政府承认的撤销}。
                  \item 事实上的承认：与对方交往又不构成法律的承认。不要求掌握。
                \end{enumerate}
          \item 继承：条约财产档案和债务（一头一尾）。协议优先。
                \begin{enumerate}
                  \item 条约：\zd{领土}划界有关的必须。
                  \item 财产档案：\qd{领土实际生存原则}，不看在哪看\qd{土地归了谁}。考不着。
                  \item 严格意义的\zd{国家非恶债}。
                  \item 一国\zd{中央政府}。\zd{地方债、个人}、其他国际法主体（\zd{商业银行}）的债务不属于。
                  \item 平等条约产生。
                  \item 合并全部；分离分立（苏联）比例。独立（殖民地）不予继承。
                \end{enumerate}
        \end{enumerate}
  \item 联合国：
        \begin{enumerate}
          \item  联大直接领导秘书处，安理会、经社理事会、托管理事会、国际法院。
                \begin{kk}{两种人员的产生}
                  \begin{enumerate}
                    \item  秘书长由\zd{安理会推荐大会简单多数}（1/2 以上）通过。先候选。
                    \item  国际法官由\zd{大会和安理会}分别选举，均\zd{特别多数（绝对多数）}获任。平行选，\zd{9} 票。
                  \end{enumerate}
                \end{kk}
          \item  三个机构的职能。
                \begin{kk}{联大}
                  \begin{enumerate}
                    \item 管得宽。
                    \item 审议权只让位于安理会。
                    \item 效力范围并不宽，内外有别。
                    \item \zd{不是立法机构}。
                  \end{enumerate}
                \end{kk}
                \begin{kk}{安理会}
                  \begin{enumerate}
                    \item  职能和平与安全。
                    \item  5 + 10 个国家。2 年。
                    \item  对于\zd{当事国和所有成员国}具有拘束力、执行力（包括武力）。
                    \item  \zd{唯一有权}采取行动的机构。相当于在执行庭。
                    \item  \zd{一国一票}平权原则。
                    \item  理事国为当事国不得投票，例外：决议包含采取行动的内容。
                    \item  程序性：\zd{9} 个同意票。国际法院大法官。
                    \item  非程序性（实质性）：\zd{要采取行动的、推荐秘书长、吸纳新会员、终止权力 \hz 开除}。\zd{9} 个同意票 + \zd{大国一致}。
                  \end{enumerate}
                \end{kk}
        \end{enumerate}
  \item 责任。
        \begin{enumerate}
          \item 战争罪扩大到个人。
          \item 绝对责任。外空行为地面损害（全部责任），核（补充责任）。
          \item 行为归因国家。各国联合军演。支配国。叛乱国成立的新国家。
          \item 个人行为：国家元首，政府首脑，外交部长，外交使节私人行为也算。履行职务。
        \end{enumerate}
  \item 免责：很少考到。
        \begin{enumerate}
          \item 权利国同意。
          \item 针对性，适度性对抗性行为。对抗措施，自卫。
          \item 不可抗力，偶然事故。
          \item 危难 \hz 紧急情况。
        \end{enumerate}
\end{enumerate}
